% Part: normal-modal-logic
% Chapter: tableaux
% Section: rules-for-K

\documentclass[../../../include/open-logic-section]{subfiles}

\begin{document}

\olfileid{nml}{tab}{rul}

\olsection{قواعد \Log{K}}

قواعد الروابط القضية المعتادة هي قواعد ال!!{tableau}s القضية الموقعة
المعتادة نفسها، ولا يطرأ على بنيتها سوى إضافة البادئات. ففي كل حالة، تنتج
القاعدة المطبقة على !!{formula} موقعة~$\sFmla{S}{!A}[\sigma]$
!!{formula}s جديدة تحمل هي أيضًا البادئة~$\sigma$. وهذا واضح حدسيًا:
فمثلًا، إذا كان~$!A \land !B$ صادقًا عند العالم الذي تسميه~$\sigma$،
كان~$!A$ و$!B$ صادقين عند~$\sigma$؛ ولا تضيف القاعدة هاتين النتيجتين
إلا بهذه البادئة، من غير أن تحكم على صدقهما في عالم آخر. ونجمع القواعد
القضية في~\olref{tab:prop-rules}.

\begin{table}
  \[\def\arraystretch{3}\begin{array}{|c|c|}
    \hline
    \AxiomC{\sFmla{\True}{\lnot !A}[\sigma]}
    \RightLabel{\TRule{\True}{\lnot}}
    \UnaryInfC{\sFmla{\False}{!A}[\sigma]}
    \DisplayProof
    &
    \AxiomC{\sFmla{\False}{\lnot !A}[\sigma]}
    \RightLabel{\TRule{\False}{\lnot}}
    \UnaryInfC{\sFmla{\True}{!A}[\sigma]}
    \DisplayProof
    \\[1ex]
    \hline
    \AxiomC{\sFmla{\True}{!A \land !B}[\sigma]}
    \RightLabel{\TRule{\True}{\land}}
    \UnaryInfC{\sFmla{\True}{!A}[\sigma]}
    \noLine
    \UnaryInfC{\sFmla{\True}{!B}[\sigma]}
    \DisplayProof
    &
    \AxiomC{\sFmla{\False}{!A \land !B}[\sigma]}
    \RightLabel{\TRule{\False}{\land}}
    \UnaryInfC{$\sFmla{\False}{!A}[\sigma] \quad \mid \quad
      \sFmla{\False}{!B}[\sigma]$}
    \DisplayProof
    \\[2ex]
    \hline
    \AxiomC{\sFmla{\True}{!A \lor !B}[\sigma]}
    \RightLabel{\TRule{\True}{\lor}}
    \UnaryInfC{$\sFmla{\True}{!A}[\sigma] \quad \mid \quad
      \sFmla{\True}{!B}[\sigma]$}
    \DisplayProof
    &
    \AxiomC{\sFmla{\False}{!A \lor !B}[\sigma]}
    \RightLabel{\TRule{\False}{\lor}}
    \UnaryInfC{\sFmla{\False}{!A}[\sigma]}
    \noLine
    \UnaryInfC{\sFmla{\False}{!B}[\sigma]}
    \DisplayProof
    \\[2ex]
    \hline
    \AxiomC{\sFmla{\True}{!A \lif !B}[\sigma]}
    \RightLabel{\TRule{\True}{\lif}}
    \UnaryInfC{$\sFmla{\False}{!A}[\sigma] \quad \mid
      \quad \sFmla{\True}{!B}[\sigma]$}
    \DisplayProof
    &
    \AxiomC{\sFmla{\False}{!A \lif !B}[\sigma]}
    \RightLabel{\TRule{\False}{\lif}}
    \UnaryInfC{\sFmla{\True}{!A}[\sigma]}
    \noLine
    \UnaryInfC{\sFmla{\False}{!B}[\sigma]}
    \DisplayProof
    \\[2ex]
    \hline
  \end{array}\]
  \caption{قواعد ال!!{tableau} ذات البادئات للروابط القضية.}
  \ollabel{tab:prop-rules}
\end{table}

شرط الإغلاق هو نفسه في ال!!{tableau}s المعتادة، إلا أننا نشترط تطابق
البادئات أيضًا، لا ال!!{formula}s وحدها. ولذلك تكون شعبة مغلقة إذا احتوت
كلًا من
\[
\sFmla{\True}{!A}[\sigma] \quad\text{و}\quad \sFmla{\False}{!A}[\sigma]
\]
لبادئة ما~$\sigma$ و!!{formula}~$!A$.

وقواعد وضع الافتراضات هي أيضًا قواعد ال!!{tableau}s المعتادة نفسها، عدا
أننا نستعمل البادئة~$1$ دائمًا مع الافتراضات. (ولا يهم أي بادئة نستعمل ما
دامت واحدة لجميع الافتراضات.) ومن ثم نقول، مثلًا، إن
\[
!B_1, \dots, !B_n \Proves !A
\]
إذا وفقط إذا وجدت لوحة مغلقة للافتراضات:
\[
\sFmla{\True}{!B_1}[1], \dots, \sFmla{\True}{!B_n}[1],
\sFmla{\False}{!A}[1].
\]

بالنسبة إلى المؤثر\iftag{prvBox}{\iftag{prvDiamond}{ين الجهويين~$\Box$
    و}{ الجهوي~$\Box$}}{}\iftag{prvDiamond}{ الجهوي~$\Diamond$}{}،
تكون بادئة نتيجة القاعدة المطبقة على !!a{formula} ذات البادئة~$\sigma$
هي~$\sigma.n$. لكن قيم~$n$ المسموح بها تتوقف على كون العلامة~$\True$ أو
$\False$.

\iftag{prvBox}{تُطيل القاعدة~$\TRule{\True}{\Box}$ شعبة تحوي
  $\sFmla{\True}{\Box !A}[\sigma]$ بإضافة
  $\sFmla{\True}{!A}[\sigma.n]$.\iftag{prvDiamond}{ وبالمثل،
    }{}}{}\iftag{prvDiamond}{تُطيل القاعدة~$\TRule{\False}{\Diamond}$
  شعبة تحوي~$\sFmla{\False}{\Diamond !A}[\sigma]$ بإضافة
  $\sFmla{\False}{!A}[\sigma.n]$.}{}
ولا يمكن تطبيق\iftag{notprvBox,notprvDiamond}{ها}{هما} إلا إذا كانت
البادئة~$\sigma.n$ قد وردت \emph{من قبل} على الشعبة التي تُطبق فيها.
ولنسم مثل هذه البادئة «مستعملة» (على الشعبة).

\iftag{prvBox}{تُطيل القاعدة~$\TRule{\False}{\Box}$ شعبة تحوي
  $\sFmla{\False}{\Box !A}[\sigma]$ بإضافة
  $\sFmla{\False}{!A}[\sigma.n]$.\iftag{prvDiamond}{ وبالمثل،
    }{}}{}\iftag{prvDiamond}{تُطيل القاعدة~$\TRule{\True}{\Diamond}$
  شعبة تحوي~$\sFmla{\True}{\Diamond !A}[\sigma]$ بإضافة
  $\sFmla{\True}{!A}[\sigma.n]$.}{}
غير أنه لا يمكن تطبيق \iftag{notprvBox,notprvDiamond}{هذه القاعدة}{هاتين
القاعدتين} إلا إذا كانت البادئة~$\sigma.n$ \emph{لم} ترد من قبل على
الشعبة التي تُطبق فيها. ونسمي مثل هذه البادئات «جديدة» (على الشعبة).

ترد القواعد في~\olref{tab:rules-K}.

\begin{table}
  \begin{center}
    \def\arraystretch{3}\def\fCenter{}
    \begin{tabular}{|c|c|}
    \hline
    \iftag{prvBox}{
      \AxiomC{\sFmla{\True}{\Box !A}[\sigma]}
      \RightLabel{\TRule{\True}{\Box}}
      \UnaryInfC{\sFmla{\True}{!A}[\sigma.n]}
      \DisplayProof
      &
      \AxiomC{\sFmla{\False}{\Box !A}[\sigma]}
      \RightLabel{\TRule{\False}{\Box}}
      \UnaryInfC{\sFmla{\False}{!A}[\sigma.n]}
      \DisplayProof\\
      $\sigma.n$ مستعملة & $\sigma.n$ جديدة
      \\[1ex]
      \hline}{}
    \iftag{prvDiamond}{
      \AxiomC{\sFmla{\True}{\Diamond !A}[\sigma]}
      \RightLabel{\TRule{\True}{\Diamond}}
      \UnaryInfC{\sFmla{\True}{!A}[\sigma.n]}
      \DisplayProof
      &
      \AxiomC{\sFmla{\False}{\Diamond !A}[\sigma]}
      \RightLabel{\TRule{\False}{\Diamond}}
      \UnaryInfC{\sFmla{\False}{!A}[\sigma.n]}
      \DisplayProof\\
      $\sigma.n$ جديدة & $\sigma.n$ مستعملة
      \\[1ex]
      \hline}{}
    \end{tabular}
  \end{center}
  \caption{القواعد الجهوية لـ~\Log{K}.}
  \ollabel{tab:rules-K}
\end{table}

إن اشتراط كون البادئة في
\iftag{prvBox}{\TRule{\True}{\Box}}{\TRule{\False}{\Diamond}}
مستعملة ضروري؛ وإلا لعددنا ما يأتي !!{tableau} مغلقة:
\iftag{notprvBox,notprvDiamond}{%
  \iftag{prvBox}{%
  \begin{oltableau}
    [\pFmla{\True}{\Box \formula{A}}{1}, just = \TAss
      [\pFmla{\False}{\lnot\Box\lnot \formula{A}}{1}, just = \TAss
        [\pFmla{\True}{\formula{A}}{1.1}, just = {\TRule{\True}{\Box}[1]}
          [\pFmla{\True}{\Box\lnot\formula{A}}{1},
            just ={\TRule{\False}{\lnot}[2]}
            [\pFmla{\True}{\lnot\formula{A}}{1.1},
              just = {\TRule{\True}{\Box}[4]}
              [\pFmla{\False}{\formula{A}}{1.1},
                  just ={\TRule{\True}{\lnot}[5]}, close]
            ]
          ]
        ]
      ]
    ]
  \end{oltableau}
  }{
  \begin{oltableau}
    [\pFmla{\True}{\lnot\Diamond\lnot \formula{A}}{1}, just = \TAss
      [\pFmla{\False}{\Diamond\formula{A}}{1}, just = \TAss
        [\pFmla{\False}{\formula{A}}{1.1}, just = {\TRule{\False}{\Diamond}[2]}
          [\pFmla{\False}{\Diamond\lnot\formula{A}}{1},
            just ={\TRule{\True}{\lnot}[1]}
            [\pFmla{\False}{\lnot\formula{A}}{1.1},
              just = {\TRule{\False}{\Diamond}[4]}
              [\pFmla{\True}{\formula{A}}{1.1},
                  just ={\TRule{\True}{\lnot}[5]}, close]
            ]
          ]
        ]
      ]
    ]
  \end{oltableau}
}}{
  \begin{oltableau}
    [\pFmla{\True}{\Box \formula{A}}{1}, just = \TAss
      [\pFmla{\False}{\Diamond \formula{A}}{1}, just = \TAss
        [\pFmla{\True}{\formula{A}}{1.1}, just = {\TRule{\True}{\Box}[1]}
          [\pFmla{\False}{\formula{A}}{1.1},
            just = {\TRule{\False}{\Diamond}[2]}, close]
        ]
      ]
    ]
\end{oltableau}}

لكن~$\Box \formula{A} \Entails/ \Diamond \formula{A}$، ولذلك سيكون
نظام برهاننا غير سليم. وكذلك~$\Diamond \formula{A} \Entails/
\Box \formula{A}$، لكن من دون شرط كون بادئة
\iftag{prvBox}{\TRule{\False}{\Box}}{\TRule{\True}{\Diamond}}
جديدة، ستكون هذه اللوحة مغلقة: \iftag{notprvBox,notprvDiamond}{%
  \iftag{prvBox}{%
  \begin{oltableau}
    [\pFmla{\True}{\lnot\Box\lnot \formula{A}}{1}, just = \TAss
      [\pFmla{\False}{\Box \formula{A}}{1}, just = \TAss
        [\pFmla{\False}{\formula{A}}{1.1}, just = {\TRule{\True}{\Box}[2]}
          [\pFmla{\False}{\Box\lnot\formula{A}}{1},
            just ={\TRule{\True}{\lnot}[1]}
            [\pFmla{\False}{\lnot\formula{A}}{1.1},
              just = {\TRule{\False}{\Box}[4]}
              [\pFmla{\True}{\formula{A}}{1.1},
                  just ={\TRule{\False}{\lnot}[5]}, close]
            ]
          ]
        ]
      ]
    ]
  \end{oltableau}
  }{
  \begin{oltableau}
    [\pFmla{\True}{\Diamond \formula{A}}{1}, just = \TAss
      [\pFmla{\False}{\lnot\Diamond\lnot\formula{A}}{1}, just = \TAss
        [\pFmla{\True}{\formula{A}}{1.1}, just = {\TRule{\True}{\Diamond}[1]}
          [\pFmla{\True}{\Diamond\lnot\formula{A}}{1},
            just ={\TRule{\False}{\lnot}[2]}
            [\pFmla{\True}{\lnot\formula{A}}{1.1},
              just = {\TRule{\True}{\Diamond}[4]}
              [\pFmla{\False}{\formula{A}}{1.1},
                  just ={\TRule{\True}{\lnot}[5]}, close]
            ]
          ]
        ]
      ]
    ]
  \end{oltableau}
}}{
  \begin{oltableau}
    [\pFmla{\True}{\Diamond \formula{A}}{1}, just = \TAss,name=one
      [\pFmla{\False}{\Box \formula{A}}{1}, just = \TAss, name=two
        [\pFmla{\True}{\formula{A}}{1.1}, just = {\TRule{\True}{\Diamond}[1]}
          [\pFmla{\False}{\formula{A}}{1.1},
            just = {\TRule{\False}{\Box}[2]}, close]
        ]
      ]
    ]
  \end{oltableau}}

\end{document}
